\documentclass[12pt,spanish,letterpaper,oneside]{article}
\def\documenttitle {Antecedentes sociales y contexto pol\'itico de los derechos humanos en M\'exico}
\def\documentsubtitle {Continuidades, transformaciones y l\'imites}
\def\documentsubject {Cuadro comparativo}
\def\documentauthor {Martin Jonathan de la Cruz}
\def\coursename {Antecedentes de los derechos humanos}
\def\coursecode {005}
\def\universityname {Universidad Abierta y a Distancia de M\'exico}
\def\universityfaculty {Licenciatura en Derecho}
\def\universitydepartment {Antecedentes de los derechos humanos}
\def\universitydepartmentimage {departamentos/UnADM}
\def\universitydepartmentimagecfg {height=1.57cm}
\def\universitylocation {Ciudad de M\'exico}
\def\authortable {\begin{tabular}{ll}
	\textbf{Alumno:} & Martin Jonathan de la Cruz \\
	Matricula: & ES2611202040 \\
	Figura docente: & Daniel Alexis Lozano Ortega \\
	Semestre/Bloque: & 2 / 2 \\
	Actividad oficial: & 5 \\
	Semana: & 4 \\
	Fecha: & 22 de octubre de 2026
\end{tabular}}
\usepackage{iftex}
\ifLuaTeX
	\let\pdfpagewidth\pagewidth
	\let\pdfpageheight\pageheight
\fi
\input{template}
\usepackage{fontspec}
\setmainfont{arial.ttf}[Path=C:/Windows/Fonts/,BoldFont=arialbd.ttf,ItalicFont=ariali.ttf,BoldItalicFont=arialbi.ttf]
\setsansfont{arial.ttf}[Path=C:/Windows/Fonts/,BoldFont=arialbd.ttf,ItalicFont=ariali.ttf,BoldItalicFont=arialbi.ttf]
\renewcommand{\familydefault}{\sfdefault}
\usepackage{pdflscape}
\usepackage{array,tabularx,colortbl}
\usepackage{tikz}
\usepackage{eso-pic}
\setcitestyle{authoryear,open={(},close={)}}
\newcommand{\insertcoverwatermark}{%
	\AddToShipoutPictureBG*{%
		\begin{tikzpicture}[remember picture,overlay]
			\node[opacity=0.16,inner sep=0pt] at (current page.center)
				{\includegraphics[width=0.72\paperwidth]{img/departamentos/UnADM.pdf}};
		\end{tikzpicture}%
	}%
}
\definecolor{cabecera}{RGB}{225,237,232}
\renewcommand{\tabularxcolumn}[1]{p{#1}}
\newcolumntype{Y}{>{\raggedright\arraybackslash}X}
\setlength{\tabcolsep}{5pt}
\setlength{\emergencystretch}{3em}
\onehalfspacing
\newcommand{\encabezado}{\rowcolor{cabecera}\textbf{Contexto hist\'orico}&\textbf{Antecedentes sociales}&\textbf{Contexto pol\'itico}&\textbf{Acontecimientos}&\textbf{Derechos o principios}&\textbf{Aportaciones y l\'imites}\\\hline}
\begin{document}
\insertcoverwatermark
\templatePortrait
\templatePagecfg
\templateFinalcfg
\newgeometry{margin=2.5cm}
\fontsize{12}{14.5}\selectfont
\onehalfspacing

\section*{Introducci\'on}
El desarrollo de los derechos humanos en M\'exico comprende conflictos
sociales, transformaciones del poder y cambios en las formas de protecci\'on
jur\'idica. No constituye una secuencia de avances inevitables: el
reconocimiento de derechos puede coexistir con exclusiones y dificultades
para ejercerlos. La historia de su concepci\'on moderna muestra avances y
retrocesos, mientras que la organizaci\'on constitucional permite distinguir
la proclamaci\'on de libertades de las instituciones que limitan el poder
\citep[p.~11]{alvarez2009}\citep[pp.~1--2]{arteaga}.

La comparaci\'on siguiente recorre ocho etapas: antecedentes prehisp\'anicos,
Conquista, Virreinato, Ilustraci\'on, Independencia, constitucionalismo liberal,
Revoluci\'on y evoluci\'on contempor\'anea. En las primeras etapas se examinan
formas de organizaci\'on y controversias sobre la protecci\'on de las personas,
sin atribuirles el concepto contempor\'aneo de derechos universales. Las
declaraciones liberales, las reivindicaciones sociales y las garant\'ias actuales
se comparan atendiendo tanto a sus aportaciones como a sus l\'imites.

\clearpage
\begin{landscape}
\pagestyle{empty}
\singlespacing
\setlength{\parindent}{0pt}
\setlength{\parskip}{0pt}
\setlength{\tabcolsep}{4pt}
\setlength{\extrarowheight}{0pt}
\renewcommand{\arraystretch}{1}
\fontsize{12}{14.5}\selectfont
\noindent\textbf{Antecedentes sociales y contexto pol\'itico: etapas 1 y 2}\par\medskip
\begin{tabularx}{\linewidth}{|Y|Y|Y|Y|Y|Y|}
\hline\encabezado
\textbf{1. Antecedentes prehisp\'anicos}\par Antes de la Conquista; sociedades diversas, no un orden jur\'idico uniforme.
& Entre los mexicas, el calpulli agrupaba familias y actividades productivas. Exist\'ian jerarqu\'ias y educaci\'on diferenciada por clase social.
& El hueytlatoani gobernaba junto al tlatocan. Los pueblos tributarios pod\'ian conservar autoridades e idiomas; el dominio no supon\'ia uniformidad.
& Organizaci\'on comunitaria y tributaria; conflictos entre pueblos sometidos y autoridades dominantes. La muestra mexica no representa a toda Mesoam\'erica.
& Responsabilidad comunitaria y orden social como referentes hist\'oricos; no equivalen a igualdad universal ni a derechos constitucionales actuales.
& Documenta organizaci\'on y jerarqu\'ias, no igualdad universal \citep[cap.~1]{brom2017}.\par La concepci\'on moderna de dignidad debe distinguirse \citep[p.~11]{alvarez2009}.\\\hline
\textbf{2. Conquista y defensa de la dignidad ind\'igena}\par Siglo XVI; encuentro conflictivo y establecimiento del dominio espa\~nol.
& La encomienda exig\'ia tributo y servicios. El deber del encomendero de proteger y evangelizar coexist\'ia con explotaci\'on y dominio.
& Expansi\'on mon\'arquica y evangelizaci\'on. La Corona conced\'ia encomiendas, pero tambi\'en las limitaba por intereses fiscales y pol\'iticos.
& Conquista de Tenochtitlan y expansi\'on territorial; acci\'on de evangelizadores. Brom destaca a Las Casas y Quiroga en la defensa de los ind\'igenas.
& Defensa de la condici\'on humana y cr\'itica del sometimiento; argumentos de su \'epoca, distintos del sistema internacional contempor\'aneo.
& La encomienda no equival\'ia a propiedad de tierras ni a protecci\'on efectiva. Sus obligaciones y abusos muestran la tensi\'on entre norma y dominio \citep[caps.~4--5]{brom2017}.\\\hline
\end{tabularx}
\clearpage
\noindent\textbf{Antecedentes sociales y contexto pol\'itico: etapas 3 y 4}\par\medskip
\begin{tabularx}{\linewidth}{|Y|Y|Y|Y|Y|Y|}
\hline\encabezado
\textbf{3. Virreinato de Nueva Espa\~na}\par Orden colonial desde el siglo XVI hasta comienzos del XIX.
& Sociedad estratificada: peninsulares, criollos, ind\'igenas y castas. Las rep\'ublicas de indios conservaron formas de tierra comunal bajo control colonial.
& Monarqu\'ia y Consejo de Indias; virrey y Audiencia. Esta administraba justicia y asesoraba y vigilaba al virrey.
& Establecimiento del r\'egimen virreinal en 1535; organizaci\'on de rep\'ublicas de indios y espa\~noles; consolidaci\'on de jerarqu\'ias coloniales.
& Legislaci\'on y jurisdicciones diferenciadas, no ciudadan\'ia igualitaria. La tutela de grupos particulares coexist\'ia con subordinaci\'on social.
& Distingue protecci\'on particular y desigualdad estructural. Tierra comunal, instituciones y jerarqu\'ias est\'an documentadas en Brom \citep[caps.~5--6]{brom2017}.\\\hline
\textbf{4. Ilustraci\'on y revoluciones liberales}\par Siglos XVII--XVIII; referentes europeos y norteamericanos.
& Cuestionamiento de privilegios, impuestos y desigualdades. Creciente disputa sobre representaci\'on y legitimidad del gobierno.
& Cr\'itica del absolutismo y afirmaci\'on de soberan\'ia, representaci\'on y separaci\'on de poderes.
& Independencia estadounidense de 1776 y Declaraci\'on francesa de 1789 como referentes de posteriores transformaciones.
& Libertad, igualdad jur\'idica, propiedad, seguridad, resistencia a la opresi\'on y limitaci\'on del poder.
& Aporta un lenguaje de derechos y legitimidad constitucional. Su proclamaci\'on no elimina de inmediato las exclusiones sociales \citep[pp.~11--13]{alvarez2009}.\\\hline
\end{tabularx}
\clearpage
\noindent\textbf{Antecedentes sociales y contexto pol\'itico: etapas 5 y 6}\par\medskip
\begin{tabularx}{\linewidth}{|Y|Y|Y|Y|Y|Y|}
\hline\encabezado
\textbf{5. Independencia}\par 1810--1821; ruptura del v\'inculo colonial y proyectos de organizaci\'on nacional.
& Desigualdad y subordinaci\'on colonial; demandas populares y conflictos de criollos con privilegios peninsulares \citep[p.~73]{delgado2010}.
& Disputa de soberan\'ia y legitimidad. La insurgencia vincul\'o independencia con representaci\'on y l\'imites al poder.
& \emph{Sentimientos de la Naci\'on} (1813) y proyecto constitucional de Apatzing\'an (1814); consumaci\'on de la independencia en 1821.
& Soberan\'ia popular, divisi\'on de poderes, abolici\'on de esclavitud y distinciones de castas; aspiraci\'on de moderar desigualdades.
& Traslada la disputa de derechos a la construcci\'on del Estado. Los proyectos y su vigencia efectiva deben distinguirse \citep[nums.~5, 12 y 15]{morelos1813}.\par\citep[p.~7]{arteaga}.\\\hline
\textbf{6. Constitucionalismo liberal del siglo XIX}\par Organizaci\'on nacional tras 1821 y Constituci\'on de 1857.
& Conflictos sobre privilegios, libertades y relaciones entre sociedad e instituciones; persistencia de desigualdad material.
& Disputa entre proyectos de Estado; consolidaci\'on del principio de supremac\'ia constitucional y de l\'imites jur\'idicos al gobierno.
& Constituciones y reformas del siglo XIX; Constituci\'on de 1857 como antecedente directo del texto reformado en 1917.
& Libertades individuales, legalidad y protecci\'on constitucional frente a actos de autoridad.
& Fortalece la organizaci\'on del poder y la defensa de libertades. La igualdad ante la ley no resuelve por s\'i sola la desigualdad econ\'omica \citep[pp.~1--2, 7]{arteaga}.\\\hline
\end{tabularx}
\clearpage
\noindent\textbf{Antecedentes sociales y contexto pol\'itico: etapas 7 y 8}\par\medskip
\begin{tabularx}{\linewidth}{|Y|Y|Y|Y|Y|Y|}
\hline\encabezado
\textbf{7. Revoluci\'on y Constituci\'on de 1917}\par Demandas revolucionarias y derechos sociales.
& Reivindicaciones agrarias, laborales y educativas; respuesta a desigualdades que no se agotan en la igualdad formal.
& Reorganizaci\'on del Estado mediante el poder constituyente; ampliaci\'on de su intervenci\'on en asuntos sociales.
& Promulgaci\'on constitucional de 1917; articulaci\'on de derechos individuales y contenidos sociales.
& Educaci\'on, r\'egimen agrario y trabajo; art\'iculos 3, 27 y 123 como referentes de protecci\'on social.
& Ampl\'ia la protecci\'on hacia condiciones materiales. El texto vigente incluye reformas: no toda su redacci\'on data de 1917 \citep[p.~2]{arteaga}.\par\citep[arts.~3, 27 y 123]{cpeum2026}.\\\hline
\textbf{8. Evoluci\'on}\par\textbf{contempor\'anea}\par Sistema p\'ublico de defensa y derechos humanos constitucionales.
& Organizaci\'on ciudadana y demandas de democracia, no discriminaci\'on, justicia y atenci\'on a violaciones de derechos.
& Apertura institucional, organismos de defensa y vinculaci\'on entre Constituci\'on y tratados internacionales.
& CNDH: 1990; rango constitucional: 1992; autonom\'ia: 1999. Reforma de derechos humanos: 2011.
& Principio pro persona; universalidad, interdependencia, indivisibilidad y progresividad; deberes de promover, respetar, proteger y garantizar.
& Articula defensa no jurisdiccional y obligaciones constitucionales. Las recomendaciones no son sentencias \citep[pp.~47--48]{alvarez2009}.\par\citep[arts.~1 y 102 B]{cpeum2026}.\\\hline
\end{tabularx}
\end{landscape}
\clearpage
\pagestyle{plain}
\onehalfspacing

\section*{Semejanzas y diferencias}
Las ocho etapas relacionan formas de organizaci\'on social con el ejercicio
del poder. Su semejanza principal es que muestran tensiones entre las
reglas proclamadas y las condiciones reales de protecci\'on. No son,
sin embargo, equivalentes: pertenecer a una comunidad o contar con una
protecci\'on particular en el orden colonial no significa ser titular de
derechos universales en el sentido constitucional actual.

La Ilustraci\'on y la Independencia comparten el cuestionamiento de la
legitimidad del poder, pero difieren en su escala: las declaraciones liberales
formulan principios generales, mientras que la insurgencia mexicana los
vincula con emancipaci\'on y construcci\'on nacional. El constitucionalismo
liberal prioriza libertades y l\'imites jur\'idicos; el social incorpora demandas
agrarias y laborales. Ambos contenidos permanecen en la Constituci\'on y
no se sustituyen entre s\'i \citep[pp.~1--2]{arteaga}.

La evoluci\'on contempor\'anea agrega organismos de defensa y una
articulaci\'on expl\'icita con tratados internacionales. Su diferencia no es
solo ampliar el cat\'alogo: el art\'iculo 1 exige obligaciones a todas las
autoridades y ordena interpretar favoreciendo la protecci\'on m\'as amplia
\citep[art.~1]{cpeum2026}. La distancia entre reconocimiento y cumplimiento
sigue siendo el punto com\'un que debe evaluarse en cada periodo.

\section*{Conclusi\'on}
El estudio comparado permite comprender que los derechos actuales no
proceden de una sola declaraci\'on ni de un progreso lineal. Sus antecedentes
muestran conflictos sobre dignidad, legitimidad del gobierno, libertades,
condiciones materiales y mecanismos de defensa. Reconocer esos cambios
impide confundir instituciones antiguas con derechos universales o atribuir
a la promulgaci\'on de una norma una eficacia autom\'atica.

Desde esta perspectiva, estudiar el pasado resulta relevante para valorar
el sistema actual: permite preguntar qui\'en recibe protecci\'on, qu\'e autoridad
est\'a obligada y qu\'e medio permite reclamarla. La importancia de un
acontecimiento no depende solo del principio que enuncia, sino tambi\'en
de las exclusiones que combate y de los instrumentos que deja disponibles.
Ese criterio ofrece una base argumentada para examinar la protecci\'on
efectiva de la dignidad humana en M\'exico.

\clearpage
\begin{thebibliography}{6}
\bibitem[\'Alvarez Icaza Longoria(2009)]{alvarez2009}
\'Alvarez Icaza Longoria, E. (2009). \emph{Los derechos humanos en M\'exico}. Nostra Ediciones.
\bibitem[Arteaga Nava(2014)]{arteaga}
Arteaga Nava, E. (2014). \emph{Derecho constitucional} (4.\textsuperscript{a} ed., edici\'on digital). Oxford University Press.
\bibitem[Brom(2017)]{brom2017}
Brom, J. (2017). \emph{Esbozo de historia de M\'exico} (1.\textsuperscript{a} ed. digital). Penguin Random House Grupo Editorial.
\bibitem[C\'amara de Diputados(2026)]{cpeum2026}
C\'amara de Diputados del H. Congreso de la Uni\'on. (2026).
\emph{Constituci\'on Pol\'itica de los Estados Unidos Mexicanos}
(\'ultima reforma DOF 7 de octubre de 2026).
\url{https://www.diputados.gob.mx/LeyesBiblio/pdf/CPEUM.pdf}
\bibitem[Delgado de Cant\'u(2010)]{delgado2010}
Delgado de Cant\'u, G. M. (2010). \emph{Historia universal: De la era de las revoluciones al mundo globalizado} (3.\textsuperscript{a} ed.). Pearson Educaci\'on.
\bibitem[Morelos y Pav\'on(1813/2021)]{morelos1813}
Morelos y Pav\'on, J. M. (2021). Sentimientos de la Naci\'on.
En M. \'A. Fern\'andez Delgado, \emph{Los Sentimientos de la Naci\'on de Jos\'e Mar\'ia Morelos}.
Instituto Nacional de Estudios Hist\'oricos de las Revoluciones de M\'exico.
(Obra original publicada en 1813).
\url{https://inehrm.gob.mx/recursos/Libros/Los_Sentimientos_de_la_Nacion_de_Jose_Maria_Morelos.pdf}
\end{thebibliography}
\noindent Asistencia de IA: GitHub Copilot apoy\'o la organizaci\'on conceptual.
\end{document}